% GENERATED FILE. Edit canonical lessons, then run npm run generate:books.
% canonical-source-hash: fnv1a64:5c0e4ad9e44656bd
% canonical-lessons: TA-C272-ataiyala-attai, TA-C272-vari, TA-C272-vinnappam, TA-C272-pukaippatam, TA-C272-nakal

\chapter{Identity card, Tax, Application, Photograph, Copy}
\label{ch:ta-identity-card-tax-application-photograph-copy}
\hlchaptermodality{272}

\begin{chapteropening}
\textbf{By the end of this chapter:} \emph{I can say an identity card, a tax, an application, a photograph and a copy.}

Complete the last lesson of chapter 272: I can say an identity card, a tax, an application, a photograph and a copy.
\end{chapteropening}

\section[\texorpdfstring{\ta{அடையாள} \ta{அட்டை} (aṭaiyāḷa aṭṭai)}{அடையாள அட்டை (aṭaiyāḷa aṭṭai)}]{\ta{அடையாள} \ta{அட்டை} (aṭaiyāḷa aṭṭai) — an identity card}

\label{lesson:TA-C272-ataiyala-attai}



\begin{quote}
Before the new one: say the Tamil for an orange, then the Tamil for a papaya.
\end{quote}

\subsection*{You'll want to know: \ta{அடையாள} \ta{அட்டை}}
\textbf{\ta{அடையாள} \ta{அட்டை}} — \emph{aṭaiyāḷa aṭṭai} — \textquotedblleft{}an identity card\textquotedblright{}.

\begin{grammarlens}[title={What you've built}]
One more word for this chapter.
\end{grammarlens}

\subsection*{Your turn}
\begin{itemize}
\raggedright
  \item \emph{Say it:} \emph{aṭaiyāḷa aṭṭai}
  \item \emph{Say it:} \emph{aṭaiyāḷa aṭṭai}, once more
  \item \emph{Say it:} say \emph{aṭaiyāḷa aṭṭai}
  \item \emph{Recall:} say the Tamil for an orange, then the Tamil for a papaya, then say \emph{aṭaiyāḷa aṭṭai} again
\end{itemize}

\begin{tcolorbox}[breakable,colback=teal!4,colframe=teal!35!black,title={Before you move on}]
What does \ta{அடையாள} \ta{அட்டை} mean? (\textquotedblleft{}An identity card\textquotedblright{}.) Say it once more.
\end{tcolorbox}
\section[\texorpdfstring{\ta{வரி} (vari)}{வரி (vari)}]{\ta{வரி} (vari) — a tax; a line (of text)}

\label{lesson:TA-C272-vari}



\begin{quote}
Before the new one: say the Tamil for a papaya, then the Tamil for an identity card.
\end{quote}

\subsection*{You'll want to know: \ta{வரி}}
\textbf{\ta{வரி}} — \emph{vari} — \textquotedblleft{}a tax; a line (of text)\textquotedblright{}.

\begin{grammarlens}[title={What you've built}]
One more word for this chapter.
\end{grammarlens}

\subsection*{Your turn}
\begin{itemize}
\raggedright
  \item \emph{Say it:} \emph{vari}
  \item \emph{Say it:} \emph{vari}, once more
  \item \emph{Say it:} say \emph{vari}
  \item \emph{Recall:} say the Tamil for a papaya, then the Tamil for an identity card, then say \emph{vari} again
\end{itemize}

\begin{tcolorbox}[breakable,colback=teal!4,colframe=teal!35!black,title={Before you move on}]
What does \ta{வரி} mean? (\textquotedblleft{}A tax; a line (of text)\textquotedblright{}.) Say it once more.
\end{tcolorbox}
\section[\texorpdfstring{\ta{விண்ணப்பம்} (viṇṇappam)}{விண்ணப்பம் (viṇṇappam)}]{\ta{விண்ணப்பம்} (viṇṇappam) — an application}

\label{lesson:TA-C272-vinnappam}



\begin{quote}
Before the new one: say the Tamil for an identity card, then the Tamil for a tax; a line.
\end{quote}

\subsection*{You'll want to know: \ta{விண்ணப்பம்}}
\textbf{\ta{விண்ணப்பம்}} — \emph{viṇṇappam} — \textquotedblleft{}an application\textquotedblright{}.

\begin{grammarlens}[title={What you've built}]
One more word for this chapter.
\end{grammarlens}

\subsection*{Your turn}
\begin{itemize}
\raggedright
  \item \emph{Say it:} \emph{viṇṇappam}
  \item \emph{Say it:} \emph{viṇṇappam}, once more
  \item \emph{Say it:} say \emph{viṇṇappam}
  \item \emph{Recall:} say the Tamil for an identity card, then the Tamil for a tax; a line, then say \emph{viṇṇappam} again
\end{itemize}

\begin{tcolorbox}[breakable,colback=teal!4,colframe=teal!35!black,title={Before you move on}]
What does \ta{விண்ணப்பம்} mean? (\textquotedblleft{}An application\textquotedblright{}.) Say it once more.
\end{tcolorbox}
\section[\texorpdfstring{\ta{புகைப்படம்} (pukaippaṭam)}{புகைப்படம் (pukaippaṭam)}]{\ta{புகைப்படம்} (pukaippaṭam) — a photograph}

\label{lesson:TA-C272-pukaippatam}



\begin{quote}
Before the new one: say the Tamil for a tax; a line, then the Tamil for an application.
\end{quote}

\subsection*{You'll want to know: \ta{புகைப்படம்}}
\textbf{\ta{புகைப்படம்}} — \emph{pukaippaṭam} — \textquotedblleft{}a photograph\textquotedblright{}.

\begin{grammarlens}[title={What you've built}]
One more word for this chapter.
\end{grammarlens}

\subsection*{Your turn}
\begin{itemize}
\raggedright
  \item \emph{Say it:} \emph{pukaippaṭam}
  \item \emph{Say it:} \emph{pukaippaṭam}, once more
  \item \emph{Say it:} say \emph{pukaippaṭam}
  \item \emph{Recall:} say the Tamil for a tax; a line, then the Tamil for an application, then say \emph{pukaippaṭam} again
\end{itemize}

\begin{tcolorbox}[breakable,colback=teal!4,colframe=teal!35!black,title={Before you move on}]
What does \ta{புகைப்படம்} mean? (\textquotedblleft{}A photograph\textquotedblright{}.) Say it once more.
\end{tcolorbox}
\section[\texorpdfstring{\ta{நகல்} (nakal)}{நகல் (nakal)}]{\ta{நகல்} (nakal) — a copy, a photocopy}

\label{lesson:TA-C272-nakal}



\begin{quote}
Before the new one: say the Tamil for an application, then the Tamil for a photograph.
\end{quote}

\subsection*{You'll want to know: \ta{நகல்}}
\textbf{\ta{நகல்}} — \emph{nakal} — \textquotedblleft{}a copy, a photocopy\textquotedblright{}.

\begin{grammarlens}[title={What you've built}]
One more word for this chapter.
\end{grammarlens}

\subsection*{Your turn}
\begin{itemize}
\raggedright
  \item \emph{Say it:} \emph{nakal}
  \item \emph{Say it:} \emph{nakal}, once more
  \item \emph{Say it:} say \emph{nakal}
  \item \emph{Recall:} say the Tamil for an application, then the Tamil for a photograph, then say \emph{nakal} again
\end{itemize}

\begin{tcolorbox}[breakable,colback=teal!4,colframe=teal!35!black,title={Before you move on}]
What does \ta{நகல்} mean? (\textquotedblleft{}A copy, a photocopy\textquotedblright{}.) Say it once more.
\end{tcolorbox}
